% HD-N13. GPT-6 Astra Pro, /root. 21 November 2026.
\section{The direct-sum replacement: an obstruction without simplicity}

\paragraph{HD-N13.}
Let $X$ be a complex abelian fourfold. No finite direct sum of
shifts of line bundles on $X$ has both the ample imaginary
two-exponential character (N11.1) and weak semiregularity for an
actual coherent action of a faithful finite point subgroup
$G\subset X\times\widehat X$.
This assertion permits arbitrary invariant endomorphism algebras
and arbitrary integral shifts; it does not assume simplicity or
vanishing of negative self-Ext groups. The proof consists of the
actual invariant direct-summand maps and the explicit real
divisor identities below.

Write the actual object as
\[
 F=\bigoplus_{j=1}^m L_j[a_j],
 \qquad L_j\in\operatorname{Pic}(X),\quad a_j\in\mathbb Z.
                                                               \tag{N13.1}
\]
The notation allows repetitions. Assume it carries the stated
coherent action and the character (N11.1); we prove that its
degree-two trace has a nonzero invariant kernel.

\paragraph{Actual invariant summands and their scalar directions.}
The category of perfect complexes on $X$ has finite-dimensional
Hom spaces and splits idempotents. Each $L_j[a_j]$ has
endomorphism algebra $\mathbb C$, so the displayed decomposition
is a Krull--Schmidt decomposition after repeated isomorphic
terms are grouped. Translation and tensoring by a flat line
bundle preserve the shift, rank and rational first Chern class.
The action therefore permutes the finitely many isomorphism
classes of displayed summands, and every orbit has a common
shift and a common first Chern class.

Here is a precise construction of invariant orbit summands even
when the given linearizations are not diagonal. Put
$E=\operatorname{End}(F)$ and let $J$ be its Jacobson radical.
Then $J$ is nilpotent, and
$E/J$ is the product of matrix algebras indexed by the distinct
isomorphism classes of line-bundle shifts in (N13.1).
The sum of the central idempotents belonging to an orbit is
$G$-invariant. Averaging gives a surjection
\[
 E^G\longrightarrow(E/J)^G.
\]
Idempotents lift through the nilpotent ideal $J\cap E^G$.
One explicit lifting procedure starts with a lift $e$ and
replaces it by
\[
 e-(2e-1)^{-1}(e^2-e).
\]
The error is squared at each step, so the procedure terminates.
Applying it in successive complementary corners gives mutually
orthogonal invariant idempotents whose sum is one. Splitting
them in the equivariant category gives actual objects $F_1,\ldots,F_k$
and an equivariant isomorphism
\[
 F\simeq\bigoplus_{\nu=1}^kF_\nu.                         \tag{N13.2}
\]
These are precisely the underlying orbit sums. To check the
identification, two idempotents $e,f$ with the same image modulo
$J$ are intertwined by the invertible element
$u=fe+(1-f)(1-e)$: one has $u\equiv1\pmod J$ and $ue=fu$.
Thus each lifted projector has the same underlying direct
summand as the original orbit projector.

Consequently
\[
 \operatorname{ch}(F_\nu)=w_\nu e^{D_\nu},
 \qquad w_\nu\in\mathbb Z\setminus\{0\},\quad
 D_\nu\in\operatorname{NS}(X)_{\mathbb Q}.                \tag{N13.3}
\]
Indeed $w_\nu$ is the signed number of line-bundle summands in
that orbit; their common shift makes its sign constant.
For each orbit, scalar multiplication gives an invariant map
\[
 H^2(X,\mathcal O_X)\longrightarrow
 \operatorname{Ext}^2(F_\nu,F_\nu)^G,
 \qquad\beta\longmapsto\beta\operatorname{id}_{F_\nu}.
                                                               \tag{N13.4}
\]
Its degree-zero trace component is $w_\nu\beta$, so it is
injective. Since $h^2(\mathcal O_X)=6$, the independent diagonal
blocks of (N13.2) imply
\[
 \dim\operatorname{Ext}^2(F,F)^G\ge6k.                    \tag{N13.5}
\]

HD-N12 identifies weak semiregularity with injectivity of the
ambient invariant trace by the actual forgetful Hochschild
isomorphism. The invariant Serre pairing and (N11.3) make its
transpose the invariant evaluation map. Thus trace injectivity
would make that evaluation surjective. The composite has
rank12 by the explicit character calculation HD-N11. Hence
\[
 \dim\operatorname{Ext}^2(F,F)^G=12,
 \qquad k\le2.                                          \tag{N13.6}
\]
This argument uses neither the Euler characteristic nor an
assumption on the endomorphism dimension.

\paragraph{Two real divisor exponentials cannot give this character.}
It remains to exclude the character identities with one or two
terms in (N13.3). Multiplication by $e^{-B}$ replaces each
$D_\nu$ in the identity by $D_\nu-B$, still a real rational
$(1,1)$ class. This is a stated operation on cohomology; no
fractional line-bundle twist of the object is required.
Write $\alpha=a+bs$, with $a,b\in\mathbb Q$ and $s^2=-d$.
Set
\[
 r=2a,\qquad c=-2db,\qquad
 A=\alpha\bar\alpha=a^2+db^2>0.
\]
The first four factorial-weighted coefficients of the target
character are
\[
 m_0=r,\qquad m_1=c,\qquad m_2=-dr,\qquad m_3=-dc,
 \qquad dr^2+c^2=4dA>0.                                 \tag{N13.7}
\]
For one exponential $w e^D$, the rank gives $w=r\ne0$ and
the degree-two class gives $D=(c/r)H$. The degree-four class
would imply $c^2+dr^2=0$, contradicting (N13.7).

Now take two nonzero weights $w_1,w_2$, two real $(1,1)$ classes
$D_1,D_2$, and write $r=w_1+w_2$. First suppose $r\ne0$.
The first and second moment identities are
\[
 w_1D_1+w_2D_2=cH,\qquad
 w_1D_1^2+w_2D_2^2=-drH^2.
\]
Subtracting the square of the first identity divided by $r$
gives the exact equality
\[
 (D_1-D_2)^2=\kappa H^2,
 \qquad\kappa=-\frac{4dA}{w_1w_2}\ne0.                   \tag{N13.8}
\]
Represent these classes by invariant real $(1,1)$ forms.
The ample form $H$ is positive definite. Diagonalize the
Hermitian form $D_1-D_2$ relative to it, writing
\[
 H=h_1+h_2+h_3+h_4,
 \qquad D_1-D_2=u_1h_1+u_2h_2+u_3h_3+u_4h_4,
 \quad u_j\in\mathbb R.
\]
Invariant forms represent the exterior cohomology algebra of
the torus faithfully. Equation (N13.8) is therefore precisely
$u_i u_j=\kappa$ for all $i\ne j$. None of the $u_j$ is zero.
The equations with one index fixed, followed by the equation
for two other indices, give
$u_1=u_2=u_3=u_4=u$ and $\kappa=u^2>0$.
Thus $D_1-D_2=uH$, and the first moment gives
$D_1=q_1H$, $D_2=q_2H$ for real numbers $q_1,q_2$.

The moments of these two real scalar exponentials satisfy
\[
 m_{j+2}-S m_{j+1}+P m_j=0,\qquad
 S=q_1+q_2,\quad P=q_1q_2.
\]
Using $j=0,1$ and (N13.7) gives
\[
 (P-d)r-Sc=0,\qquad (P-d)c+Sdr=0.
\]
Their coefficient determinant is $dr^2+c^2=4dA>0$.
Consequently $S=0$ and $P=d$. But then
$P=q_1(-q_1)=-q_1^2\le0$, contradicting $d>0$.

Finally suppose $r=0$. Then $w_2=-w_1$, $c\ne0$, and the
first moment gives $D_1-D_2=(c/w_1)H$. The second moment is
zero, so
\[
 H(D_1+D_2)=0.
\]
This implies $D_1+D_2=0$. One direct verification diagonalizes
that real Hermitian form relative to $H$; its eigenvalues
$v_j$ satisfy $v_i+v_j=0$ for every pair, forcing all four to
vanish. Hence $D_1=uH$, $D_2=-uH$, where $2w_1u=c\ne0$.
The third moment is $cu^2H^3$, whereas (N13.7) requires
$-dcH^3$. Thus $u^2=-d$, again impossible for real $u$.

These actual coefficient and Hermitian-form calculations
exclude every one- or two-orbit character identity. They
contradict (N13.6), completing the proof. \hfill$\square$

\paragraph{Scope.}
This result treats actual direct sums of shifted line bundles,
including coherent point-group actions that mix their chosen
summands. It leaves nontrivial extensions, complexes with other
indecomposable constituents, and direct constructions on an
eightfold as separate questions. The obstruction is the exhibited
trace map and real-divisor moment calculation; it does not
identify a virtual line-bundle expansion with an actual
semiregular object.
